\section*{Záver}
\addcontentsline{toc}{section}{Záver}



Cieľom tejto diplomovej práce bolo navrhnúť implementáciu riadiaceho systému pre vybrané laboratórne zariadenie predstavujúce riadený systém.
Hlavný dôraz bol kladený na praktickú realizáciu riadiaceho systému v rámci dostupných
implementačných platforiem, konkrétne v prostredí MATLAB/Simulink, v jazyku
C++ na vývojovej doske Arduino UNO a na priemyselnom PLC systéme Siemens
SIMATIC S7-1200 G2.

V rámci práce bol najskôr spracovaný opis použitého laboratórneho zariadenia TS05,
spôsob jeho pripojenia k jednotlivým platformám a základné postupy potrebné
pre prácu s jeho vstupnými a výstupnými signálmi. Pre potreby návrhu regulátorov
bola vykonaná identifikácia systému na základe experimentálne nameraných údajov.
Získaný model bol použitý pri návrhu regulačných algoritmov a pri ich počiatočnom
simulačnom overení.

Hlavným výsledkom práce je realizácia troch algoritmov riadenia na reálnom
laboratórnom zariadení. Prvým z nich je PI regulátor, ktorý predstavuje
základné riešenie pre riadenie v okolí jedného pracovného bodu. Druhým je riadiaci systém, ktorý je implementáciou vlastného návrhu PI regulácie s prepínaním pracovných bodov, ktorý
rozšíril použiteľný pracovný rozsah riadenia kombináciou globálnej a lokálnej
regulačnej vetvy. Tretím bol algoritmus z oblasti adaptívneho riadenia s referenčným modelom (MRAC) založený na gradientnej
metóde, ktorý umožnil overiť implementáciu riadiaceho systému z kategórie pokročilých metód automatického riadenia na reálnom hardvéri.

Navrhnuté algoritmy boli implementované na všetkých uvažovaných platformách.
Tým sa podarilo overiť nielen ich regulačný princíp, ale aj praktické aspekty
spojené s realizáciou riadiaceho systému, ako sú časovanie výpočtu, uchovávanie
stavových premenných, škálovanie analógových signálov, saturácia akčného zásahu
a časový záznam signálov. Porovnanie výsledkov medzi platformami
slúžilo ako doplnkové vyhodnotenie a ukázalo, že pri jednoduchších regulačných
štruktúrach možno dosiahnuť veľmi podobné správanie, zatiaľ čo adaptívny
algoritmus je výraznejšie ovplyvnený vlastnosťami konkrétnej implementácie.

Z experimentálnych výsledkov vyplynulo, že PI regulátor je vhodný najmä pre
lokálne riadenie v okolí pracovného bodu. Algoritmus s prepínaním pracovných
bodov umožnil realizovať prechody medzi rôznymi pracovnými oblasťami systému
a následne pokračovať v lokálnej regulácii. Pri MRAC regulátore sa potvrdila
možnosť implementácie adaptačného mechanizmu, zároveň sa však prejavila jeho
citlivosť na prítomnosť šumu v spätnoväzbovom signále, numerickú deriváciu, drift pracovného bodu a~voľbu
adaptačných koeficientov. Heuristická úprava použitá pri C++ implementácii
znížila drift parametrov, ale úplne ho neodstránila.

Možno konštatovať, že stanovené ciele práce boli splnené. Bol zostavený riadiaci
systém pre reálny laboratórny dynamický systém, pričom bola navrhnutá implementácia niekoľkých algoritmov riadenia a ich funkčnosť bola overená experimentálne. Práca
zároveň ukázala, že prenos algoritmu zo simulačného prostredia na reálny hardvér
nie je iba formálnym prepísaním kódu, ale samostatnou technickou úlohou, ktorá
vyžaduje zohľadnenie vlastností merania, akčných členov, výpočtovej platformy
a~fyzikálneho správania riadeného systému.

Práca tiež vytvára základ pre modifikáciu a~nadviazanie na prezentované výsledky napríklad z~hľadiska zlepšenia prezentovaného adaptívneho algoritmu riadenia v~zmysle jeho robustnosti, teda overiť normalizáciu
regresora, projekčný algoritmus, \mbox{$\sigma$-modifikáciu} alebo iné stabilizačné úpravy
zákona adaptácie. Pri algoritme s~prepínaním pracovných bodov by bolo možné
rozšíriť počet pracovných bodov, doplniť interpoláciu parametrov lokálnych
regulátorov alebo upraviť podmienku detekcie ustálenia. Samostatným smerom
ďalšej práce by mohlo byť rozšírenie návrhu zo SISO prístupu na využitie MIMO charakteru
systému TS05.